\documentclass[11pt]{article}
\usepackage{preamble}
\renewcommand{\answer}[1]{}

\begin{document}

\ladderheading{AP Statistics \textperiodcentered\ Unit 2 \textperiodcentered\ Topic 2.12 \textperiodcentered\ B: 2026-11-13 \textperiodcentered\ E: 2026-12-09}{When Is the Mean Nearly Normal?}

\focusbanner{TODAY'S PROBLEM}

Individual commute times in a very large city are strongly right-skewed, with population mean $\mu=25$ minutes and standard deviation $\sigma=15$ minutes. Consider independent random samples.\par\smallskip \textbf{Rae:} ``Any average is approximately normal, even for $n=4$.''\par \textbf{Sol:} ``With this long right tail, $n=4$ may be too small; $n=40$ is more defensible.''\par
\begin{center}
\begin{tikzpicture}
\begin{axis}[
  width=0.77\linewidth,
  height=2.07in,
  ybar interval, ymin=0, ymax=32,
  xmin=0, xmax=80, xtick={0,10,20,30,40,50,60,70,80},
  xlabel={Individual commute time (minutes)}, ylabel={Relative frequency (percent)},
  ymajorgrids, tick label style={font=\small}, label style={font=\small},
  bar shift=0pt, x tick label as interval=false
]
\addplot[fill=calloutblue, draw=dayblue] coordinates {
  (0,12)
  (10,30)
  (20,28)
  (30,15)
  (40,8)
  (50,4)
  (60,2)
  (70,1)
  (80,0)
};
\end{axis}
\end{tikzpicture}
\end{center}
\begin{firsttakebox}{FIRST TAKE --- before the video (5 min)}
Which would you expect to smooth out the long commute times more: averaging 4 people or 40? Give one reason.
\workspace{0.9in}
\end{firsttakebox}

\begin{callout}{\IconBook\ CLT: SHAPE, SIZE, AND THE RIGHT DISTRIBUTION (keep on the page)}{calloutgreen}
The sampling distribution of $\bar{x}$ contains the sample means from all possible random samples of one size $n$.
For independent observations, $\mu_{\bar{x}}=\mu$ and $\sigma_{\bar{x}}=\dfrac{\sigma}{\sqrt{n}}$.
The \textbf{central limit theorem} says that when $n$ is sufficiently large, the sampling distribution of
$\bar{x}$ is approximately normal even if the population is not. Stronger population skew generally requires
a larger $n$; a small sample can retain that skew. The CLT concerns sample means, not individual observations.
\end{callout}

\newpage
\textbf{Finish after the video:}\par\smallskip

\dokbadge{1}~\textbf{(a)} Identify the population shape and state the population mean and standard deviation with units.\par
\workspace{0.7in}

\dokbadge{2}~\textbf{(b)} For $n=4$ and $n=40$, calculate the mean and standard deviation of the sampling distribution of $\bar{x}$.\par
\workspace{1.4in}

\dokbadge{3}~\textbf{(c)} In 3--4 sentences, decide whose claim is better supported. Use the population shape, independence, and the two sample sizes. Explain which distribution the CLT describes and why using a normal curve for averages of 4 could be unreliable.\par
\workspace{2.0in}

\begin{sentenceframebox}\raggedright\textbf{Frame:} For 4 people, \blank[2.5in]; for 40 people, \blank[2.5in].\end{sentenceframebox}

\begin{sentenceframebox}\raggedright\textbf{Frame:} The CLT describes \blank[2.2in], so using a normal curve for 4 could \blank[2.2in].\end{sentenceframebox}

\begin{summaryexitbox}{EXIT --- turn this in}
One thing the video changed about my first take: \hrulefill\par
\workspace{0.4in}
\end{summaryexitbox}

\end{document}
